The growing demand for processing in vehicular applications has driven the adoption of \gls{VEC} systems. In this paradigm, vehicles with restricted computational capacity can overcome their local limitations through task offloading to edge servers and neighboring vehicles. However, optimizing processing decisions under strict deadline and energy consumption constraints is a core challenge due to the non-stationarity of these environments. The continuous change in the number of candidate nodes and the variation in workload hinder the direct use of traditional machine learning models, which often require fixed-dimension input architectures and constant retraining in the face of these variations. To overcome this limitation, this work proposes and evaluates a hybrid decision architecture based on the integration of multicriteria methods with \gls{DRL}. The multicriteria component continuously weighs the most critical aspects of decision-making, such as the estimated completion time, energy cost, and distance, providing a compact representation of the state observed by the agent while the \gls{DRL} learns the best long-term offloading policy. This hybrid solution guarantees scale invariance and learning stability, allowing the intelligent agent to operate without policy collapse regardless of the vehicular density. Evaluations in simulation environments demonstrate that the proposed architecture achieves an operational deadline compliance rate of 45.21\%, significantly outperforming multiple reference heuristics and other \gls{DRL}-based strategies. These results attest to the systemic efficiency of the model for complex scenarios. As an extension of this proposal, fuzzy logic (\textit{Fuzzy}) will be incorporated to the decision criteria, with the aim of improving the resilience of the solution against imprecise information.

\keywords{Task Offloading. Vehicular Networks. Vehicular Edge Computing. Vehicular Communication. Deep Reinforcement Learning. Fuzzy-Topsis. Multicriteria Decision.}
